\chapter{критический аудит-допуск полярного коннектора Мориты}
\label{ch:v8-polar-morita-connector-admission-gate}

\section{Проверяемая зацепка}

Том~VII оставил факторную развилку
\begin{equation}
 M_{22}(\mathbb C)\oplus M_{21}(\mathbb C)
\end{equation}
и полный угол $\operatorname{Hom}(\mathbb C^{21},\mathbb C^{22})$
размерности $462$. Алгебраически один ненулевой элемент этого угла вместе
с полными диагональными факторами порождает $M_{43}(\mathbb C)$. Поэтому
возник кандидат
\begin{equation}
 T_0=\operatorname{diag}(I_{11},U^*),
 \label{eq:v8-polar-formal-connector}
\end{equation}
где $U:\mathbb C^{11}\to\mathbb C^{10}$ есть уже выведенная полярная
коизометрия фоновой инцидентности.

Первичная проверка матриц дала
\begin{equation}
 T_0^*T_0=I_{21},\qquad
 T_0T_0^*=I_{11}\oplus P,
 \qquad I_{11}-P=Q,\quad\rank Q=1.
\end{equation}
Но перед физическим допуском необходимо проверить смысл двух совпавших
одиннадцатимерных пространств.

\section{Два разных значения числа одиннадцать}

В рёберном Ходж-комплексе
\begin{equation}
 \mathcal K_E=\mathbb C^{E_{\rm new}}\oplus
              \mathbb C^{E_{\rm new}},
 \qquad |E_{\rm new}|=11.
 \label{eq:v8-polar-edge-label-space}
\end{equation}
Здесь базис первого $\mathbb C^{11}$ состоит из одиннадцати \emph{меток
разрешённых стрелок}. Это сжатие кратностей полного модуля
\begin{equation}
 \mathcal E_{\rm new}
 =\bigoplus_{e\in E_{\rm new}}
   \operatorname{Hom}(\mathcal H_{s(e)},\mathcal H_{t(e)}),
 \qquad \dim_{\mathbb C}\mathcal E_{\rm new}=36.
 \label{eq:v8-polar-full-arrow-module}
\end{equation}

В связывающем концевом носителе
\begin{equation}
 \mathcal K_V=\mathbb C^{11}_{\rm state}\oplus\mathbb C^{10}_{\rm state},
\end{equation}
первое пространство есть физический левый модуль
\begin{equation}
 \mathbb C^{11}_{\rm state}
 =Q_L^{(6)}\oplus L_L^{(2)}\oplus X_L^{(1)}\oplus Y_L^{(2)}.
 \label{eq:v8-polar-left-state-space}
\end{equation}
Следовательно, равенство размерностей не задаёт оператор
$I_{11}:\mathbb C^{11}_{\rm state}\to\mathbb C^{E_{\rm new}}$.
Формула~\eqref{eq:v8-polar-formal-connector} скрыто выбрала такой
изоморфизм.

\section{Неустранимая свобода отождествления}

После явного обозначения скрытого выбора формальный кандидат имеет вид
\begin{equation}
 T_j=\operatorname{diag}(j,jU^*),
 \qquad
 j:\mathbb C^{11}_{\rm state}\longrightarrow
       \mathbb C^{E_{\rm new}}.
 \label{eq:v8-polar-j-family}
\end{equation}
Условие изометрии требует $j\in U(11)$, но не выбирает его. Пространство
кандидатов имеет вещественную размерность
\begin{equation}
 \dim_{\mathbb R}U(11)=121.
\end{equation}

Пусть $X_E$ есть обмен двух копий пространства меток стрелок, а
$X_V=\left(\begin{smallmatrix}0&U^*\\U&0\end{smallmatrix}\right)$ ---
полярный обмен физических пространств концевых секторов. Для каждого $j\in U(11)$
получаем
\begin{equation}
 X_ET_j-T_jX_V
 =\begin{pmatrix}0&0\\jQ&0\end{pmatrix},
 \qquad
 \|X_ET_j-T_jX_V\|_{\rm HS}=1.
 \label{eq:v8-polar-j-independent-defect}
\end{equation}
Поэтому Ходж-дефект видит индексную линию, но совершенно не различает
$121$ параметр отождествления $j$. Восемь случайных унитарных $j$ дали одну
и ту же норму дефекта, оставаясь попарно далеко друг от друга.

Сохранение линии $\ker A_0$ действительно дополняет $U^*$ до унитарной
карты между физическими пространствами состояний размерности $11$. Однако оно
не создаёт отсутствующее отображение из пространства состояний в пространство
меток стрелок. Поэтому индексное дополнение не исправляет типизацию
$I_{11}$.

\section{Что остаётся верным алгебраически}

Если произвольно выбрать $j$, то $T_j$ является ненулевым элементом
$M_{22\times21}(\mathbb C)$. Для любого его ненулевого элемента
\begin{equation}
 E_{ap}T_jE_{qb}=(T_j)_{pq}E_{ab},
\end{equation}
поэтому лево-правый линейная оболочка имеет ранг $462$, а коммутант
$M_{22}\oplus M_{21}$ уменьшается с двух скаляров до одного. Условно
получается
\begin{equation}
 \operatorname{Alg}^*(M_{22}\oplus M_{21},T_j,T_j^*)=M_{43}(\mathbb C).
\end{equation}
Это уточняет лишь арифметику: $462$ есть размер порождаемого бимодуля, а не
обязательное число независимых коэффициентов. Но вместо них остаётся
невыведенный физический изоморфизм $j$. Связывающая алгебра не выбирает
этот элемент сама \cite{v8-polar-morita-bert}; конечная спектральная тройка
требует заданных представлений и допустимого оператора Дирака
\cite{v8-polar-morita-krajewski,v8-polar-morita-cacic}.

\begin{theorem}[Запрет размерностного полярного коннектора]
Формальный оператор $\operatorname{diag}(I_{11},U^*)$ является корректной
изометрией только после невыведенного отождествления пространства
одиннадцати стрелочных меток с одиннадцатимерным физическим модулем состояний.
Существующие полярная, Ходж- и индексная структуры не выбирают это
отождествление: допустимо всё семейство $U(11)$. Поэтому первый примитив
следующей программы ещё не получен.
\end{theorem}

\begin{nogocriterion}[Запрет отождествления по размерности]
Нельзя использовать матрицу $I_{11}$ между двумя носителями только потому,
что их координатные модели имеют одинаковую размерность. Следующий маршрут
допустим лишь после классификации эквивариантных отображений от полного
физического модуля стрелок $\mathcal E_{\rm new}$ к модулю концевых секторов с
учётом представления конечной алгебры, вещественной структуры и градуировки.
\end{nogocriterion}

\begin{protocol}
\begin{enumerate}
 \item формальная изометрия после выбора $j$: \textbf{получена};
 \item размер полного орбитального угла: \textbf{$462$};
 \item $462$ независимых коэффициента обязательны: \textbf{нет};
 \item первый $\mathbb C^{11}$: \textbf{пространство меток стрелок};
 \item второй $\mathbb C^{11}$: \textbf{физический модуль состояний};
 \item полный модуль одиннадцати стрелок: \textbf{размерность $36$};
 \item канонический $I_{11}$ между носителями: \textbf{не выведен};
 \item остаточная свобода: \textbf{$U(11)$, 121 вещественный параметр};
 \item индексная линия выбирает $j$: \textbf{нет};
 \item полярный коннектор физически допущен: \textbf{нет};
 \item новый том: \textbf{пока не открывается};
 \item следующий гейт: \textbf{классификация физического эквивариантного
 интертвинера стрелки--концевой сектор}.
\end{enumerate}
\end{protocol}

Машинный сертификат сохранён в
\path{s2t_v8_polar_morita_connector_admission_gate_results.json}.

\begin{thebibliography}{9}
\bibitem{v8-polar-morita-bert}
J.~Bertozzini, R.~Conti, W.~Lewkeeratiyutkul,
\emph{A Spectral Theorem for Imprimitivity $C^*$-Bimodules},
arXiv:0812.3596.
\bibitem{v8-polar-morita-krajewski}
T.~Krajewski,
\emph{Classification of Finite Spectral Triples},
J. Geom. Phys. 28 (1998), 1--30; arXiv:hep-th/9701081.
\bibitem{v8-polar-morita-cacic}
B.~\v{C}a\'ci\'c,
\emph{Moduli Spaces of Dirac Operators for Finite Spectral Triples},
J. Noncommut. Geom. 3 (2009), 453--495; arXiv:0902.2068.
\end{thebibliography}